\section{实验结果与分析}\label{sec:results}

\subsection{六模型基准比较}
表\ref{tab:baseline-results}比较当前协议下的六模型。原XGBoost在开发与B的MSE均最低，随机森林均为第二，但两集合完整排序并不相同；MSE排名的Spearman相关为0.885714。
\begin{table}[htbp]
\small
\setlength{\tabcolsep}{5pt}
\centering
\caption{六基准在当前开发/B的误差（未舍入预测）}
\label{tab:baseline-results}
\begin{tabular}{lrrrr}
\toprule
模型 & 开发MSE & B MSE & 开发MAE & B MAE \\
\midrule
原XGBoost & 2.228808 & 1.252744 & 0.780951 & 0.648007 \\
随机森林 & 2.276036 & 1.334292 & 0.786652 & 0.664520 \\
原大网络 & 2.375530 & 1.342265 & 0.748643 & 0.596760 \\
Ridge & 2.289115 & 1.372016 & 0.807902 & 0.711438 \\
决策树 & 2.595176 & 1.401940 & 0.829650 & 0.680730 \\
AdaBoost.R2 & 2.455500 & 1.444966 & 0.808051 & 0.709821 \\
\bottomrule
\end{tabular}
\end{table}


原大网络的开发MAE为0.748643、B MAE为0.596760，分别低于原XGBoost的0.780951和0.648007。这表明同一模型对平均绝对误差与大误差平方损失的表现可能不同。因此，“原XGBoost较优”仅指当前六候选在MSE目标下的点估计，不能扩展为所有指标全面占优。后续主方案仍以MSE作训练选择目标。

\subsection{顺序优化阶段比较}
表\ref{tab:stages}展示从原XGBoost到四模型融合的阶段结果。训练CV列为两套合并OOF损失的均值，属于选择分数；开发与B均为冻结后评价。
\begin{table}[htbp]
\small
\setlength{\tabcolsep}{5pt}
\centering
\caption{顺序优化阶段的选择分数与冻结评价}
\label{tab:stages}
\begin{tabular}{L{0.31\linewidth}L{0.21\linewidth}L{0.18\linewidth}L{0.18\linewidth}}
\toprule
方案 & 两套OOF均值 & 开发MSE & B MSE \\
\midrule
原XGBoost & 1.806988 & 2.228808 & 1.252744 \\
组合特征 & 1.786509 & 2.218877 & 1.275351 \\
联合XGBoost & 1.781861 & 2.236749 & 1.242601 \\
三模型等权 & 1.762972 & 2.196540 & 1.245558 \\
去乙肝小网络 & 1.911642 & 2.260269 & 1.352561 \\
四模型各25\% & 1.759568 & 2.177243 & 1.234296 \\
\bottomrule
\end{tabular}
\end{table}


组合特征降低训练OOF与开发MSE，却使B MSE升至1.275351；联合搜索的XGBoost在B降至1.242601，开发则略差于基线。两者不支持“特征越多或搜索后必然在所有集合改善”的结论。三模型等权在本轮两个集合均低于原XGBoost，加入小网络后四模型点估计进一步降低。

小网络单独的开发/B MSE为2.260269/1.352561，并非当前最强单模型；其融合价值应由联合OOF误差检验，而非单模型排名推断。初始128--64--16加PCA网络的两套OOF均值为2.656852，明显高于去乙肝64--32网络的1.911642；本文保留该中间候选记录，不将网络结构缩小本身解释为经过完全控制的因果比较。

\subsection{最终模型的综合指标}
最终四模型开发MSE为2.177243，B MSE为1.234296。相对原XGBoost，按相同集合基线计算的下降比例为
\begin{equation}
\rho_E=\frac{\operatorname{MSE}_{E,\mathrm{baseline}}-
\operatorname{MSE}_{E,\mathrm{final}}}{\operatorname{MSE}_{E,\mathrm{baseline}}}\times100\%.
\label{eq:improvement}
\end{equation}
开发和B的$\rho_E$分别为2.313554\%和1.472594\%，不等于分类准确率增加同样百分点。表\ref{tab:final-metrics}同时给出辅助指标。
\begin{table}[htbp]
\small
\setlength{\tabcolsep}{5pt}
\centering
\caption{最终四模型的开发/B综合指标}
\label{tab:final-metrics}
\begin{tabular}{lrrrr}
\toprule
评价集合 & MSE & RMSE & MAE & $R^2$ \\
\midrule
开发 & 2.177243 & 1.475548 & 0.777364 & 0.205397 \\
B & 1.234296 & 1.110989 & 0.638160 & 0.196253 \\
\bottomrule
\end{tabular}
\end{table}


这些数值来自固定模型的已保存预测，不是竞赛平台线上分数或名次。开发评分与B评分分别使用4649与6642条拟合行，二者绝对MSE不直接作跨集合优劣判断。

\subsection{训练选择证据与互补性}
原基线的两套OOF MSE为1.803531和1.810445，三模型为1.758488和1.767456，最终四模型为1.754295和1.764841。四模型在两套划分下均低于对应三模型，构成本轮选择其加入小网络的依据。

树/线性成员训练OOF残差相关仍较高：优化XGBoost与LightGBM约0.973072，Ridge与LightGBM约0.941305。相关性较低不自动保证融合有效，最终仍以组合损失比较。已有相关矩阵未包含小网络，因此不据此补造网络与树的相关系数。

9组网络补充融合中，网络25\%加三模型75\%的两套OOF均值为1.759568；网络50\%和75\%对应均值分别为1.783212和1.833903。这是既定网格中的比较结果，不证明连续权重空间的全局最优。所有CV均参与候选选择，不能把这组差异当作新的独立泛化验证。
